\documentclass[10pt,a4paper]{article}
\usepackage{amsmath,amssymb,amsthm,graphicx}
% Minimal retained preamble for lane-p2 cards taken from
% sources/ahl-16-fixed-points-2019/editorial-preamble.tex
% (doi:10.5802/ahl.16, Annales Henri Lebesgue 2 (2019) 149-185, CC BY 4.0).
%
% Copied verbatim from the editorial preamble: the package list, the
% theorem declarations under the shared counter and the macro block, so that
% every macro and environment of the retained spans is available.
%
% Omitted and why: the class line (the assembler supplies the class); the
% a4paper geometry line and the fontspec block with its explicit lmroman10
% files (encodings and fonts fight the forced xelatex engine, and the
% assembler supplies the page setup).
\usepackage{amsmath,amssymb,amsthm,mathtools,graphicx,xcolor,hyperref,xurl,xspace,ifthen}
\newtheorem{theo}{Theorem}[section]
\newtheorem{lemm}[theo]{Lemma}
\newtheorem{theorem}[theo]{Theorem}
\newtheorem*{remas*}{Remarks}
\providecommand{\og}{«\nobreakspace}
\providecommand{\fg}{\nobreakspace»}
\newtheorem{lem}[theo]{Lemma}
\newtheorem{lemma}[theo]{Lemma}
\newtheorem{prop}[theo]{Proposition}
\newtheorem{coro}[theo]{Corollary}
\newtheorem{corol}[theo]{Corollary}
\newtheorem{conj}[theo]{Conjecture}
\newtheorem{Proposition}[theo]{Proposition}
\newtheorem{theoreme}[theo]{Theorem}
\theoremstyle{definition}
\newtheorem{defi}[theo]{Definition}
\newtheorem{exam}[theo]{Example}
\newtheorem{exem}[theo]{Example}
\newtheorem{rema}[theo]{Remark}
\newtheorem{nota}[theo]{Notation}
\newtheorem*{rema*}{Remark}
\newtheorem*{defi*}{Definition}
\newtheorem*{theo*}{Theorem}
\newtheorem*{lemm*}{Lemma}
\newtheorem*{prop*}{Proposition}
\newcommand{\enoncename}{}
\newtheorem{corpusenonce}[theo]{\enoncename}
\newenvironment{enonce}[2][]{\renewcommand{\enoncename}{#2}\begin{corpusenonce}}{\end{corpusenonce}}
\newenvironment{enonce*}[2][]{\par\medskip\noindent\textbf{#2.}\itshape}{\par\medskip}
\newenvironment{proofc}[1][\proofname]{\begin{proof}[#1]}{\end{proof}}
\newcommand{\Mc}{\mathcal{M}}
\newcommand{\pointir}{.\ ---\ }
\providecommand{\mathof}[1]{\mathrm{#1}}

\numberwithin{equation}{section}
\numberwithin{figure}{section}
\numberwithin{table}{section}
\def\endappendix{}
%~Mouliné par MaN_auto v.0.21.4 2019-04-30 09:02:09


\usepackage[all]{xy}

\hyphenation{auto-morphism
every
Roelcke
peri-od-ic
peri-od-icy
rep-re-sen-ta-tion
}








\newcommand{\A}{\mathcal A}
\newcommand{\N}{\mathbb{N}}
\newcommand{\Z}{\mathbb{Z}}
\newcommand{\Q}{\mathbb{Q}}
\newcommand{\K}{\mathcal{K}}
\newcommand{\CC}{\mathbb{C}}

\newcommand*{\mA}{\mathbf{A}}
\newcommand*{\mB}{\mathbf{B}}
\newcommand*{\mC}{\mathbf{C}}
\newcommand*{\mD}{\mathbf{D}}
\newcommand*{\mE}{\mathbf{E}}
\newcommand*{\mF}{\mathbf{F}}
\newcommand*{\mH}{\mathbf{H}}
\newcommand*{\mZ}{\mathbf{Z}}
%\newcommand*{\mA}{\mathbf{A}}
%\newcommand*{\mA}{\mathbf{A}}

\newcommand{\restrict}[2]{#1 \upharpoonright #2}

\newcommand{\e}{\varepsilon}

\DeclareMathOperator{\Aut}{Aut}
\DeclareMathOperator{\Stab}{Stab}
\DeclareMathOperator{\Age}{Age}

\newcommand{\actson}{\curvearrowright}

\newcommand{\RUCB}{\mathrm{RUC}_{b}}
\newcommand{\ROB}{\mathrm{Ro}_{b}}
\newcommand{\WAP}{\mathrm{WAP}}
\newcommand{\AP}{\mathrm{AP}}

\graphicspath{{./figures/}}

\newcommand*{\mk}{\mkern -1mu}
\newcommand*{\Mk}{\mkern -2mu}
\newcommand*{\mK}{\mkern 1mu}
\newcommand*{\MK}{\mkern 2mu}

\hypersetup{urlcolor=purple, linkcolor=blue, citecolor=red}

\newcommand*{\romanenumi}{\renewcommand*{\theenumi}{\roman{enumi}}}
\newcommand*{\Romanenumi}{\renewcommand*{\theenumi}{\Roman{enumi}}}
\newcommand*{\alphenumi}{\renewcommand*{\theenumi}{\alph{enumi}}}
\newcommand*{\Alphenumi}{\renewcommand*{\theenumi}{\Alph{enumi}}}
\let\oldexists\exists
\renewcommand*{\exists}{\mathrel{\oldexists}}
\let\oldforall\forall
\renewcommand*{\forall}{\mathrel{\oldforall}}


\usepackage{cleveref}
\crefname{originalproposition}{original proposition}{original propositions}
\Crefname{originalproposition}{Original proposition}{Original propositions}
\makeatletter
\newlabel{prop:RPExtRoelcke}{{4.9}{source}{}{sourcepointer-1}{}}
\newlabel{prop:RPExtRoelcke@cref}{{[originalproposition][0][]4.9}{[1][0][]source}{}{}{}}
\makeatother
\crefname{originalproposition}{original proposition}{original propositions}
\Crefname{originalproposition}{Original proposition}{Original propositions}
\makeatletter
\newlabel{prop:Roelcke-precompact}{{4.2}{source}{}{sourcepointer-2}{}}
\newlabel{prop:Roelcke-precompact@cref}{{[originalproposition][0][]4.2}{[1][0][]source}{}{}{}}
\makeatother
\title{The Ramsey property for the Roelcke colorings of a Fraisse structure is equivalent to the definable Ramsey property of its age}
\author{Lionel Nguyen Van Thé}
\date{doi:10.5802/ahl.16; CC BY 4.0}
\begin{document}
\maketitle
\noindent\textit{Editorial note.} The counted unit is the article's own Proposition prop:DefinableRP at source lines 580-582 together with its proof at source lines 584-588, a straight adjacent pairing in the source file. The source is the publisher version of the article in Annales Henri Lebesgue 2 (2019), pages 149-185, DOI 10.5802/ahl.16, released under CC BY 4.0.\par\medskip
\noindent\textit{Editorial note.} The licence evidence used here is the observation file kept with the source, license-observation.json, which records the licence as CC BY 4.0, points at the inherited publisher-version evidence and states that the landing page displays no licence line; the underlying evidence files are the same ones the earlier cards of this source rely on. No acceptance record is written anywhere.\par\medskip
\noindent\textit{Editorial note.} The retained proof names two earlier propositions of the article: the one on the extension of Roelcke colorings and the one characterising Roelcke precompactness. Neither is retained in this card, so both are declared as navigation pointers. Their numbers are the ones the article prints, and the article declares its theorem-like environments under a single counter attached to the section, so the printed number of each statement is the section together with its place inside that section: 4.9 and 4.2 respectively. The count was made by enumerating every numbered theorem-like environment of the source in order and reading the indices off that enumeration. Earlier revisions of this card carried 36 and 29 and then 35 and 28, all of which were place-in-the-article counts rather than printed numbers; revision r003 carries the printed numbers and was re-assembled and checked. The anchors recorded point at each statement's own lines.\par\medskip
\noindent\textit{Editorial note.} The retained preamble is the source's editorial preamble with three omissions only: the class line, which the assembler supplies; the page geometry; and the fontspec block with its explicit font files, which would fight the forced xelatex engine. Everything else, including the theorem declarations under the shared counter and the whole macro block, is retained verbatim.\par\medskip
\noindent\textit{Editorial note.} Neither the retained statement nor the retained proof contains a citation, so this card needs no bibliography entry.\par\medskip
\noindent\textit{Editorial note.} The census package marked this candidate free of missing citation keys and listed the statement span as 580-582 and the proof span as 584-588; both were re-read from the source before the card was written. Revision r003 of this card replaces the place-in-the-article numbering of the two pointers with the section-scoped printed numbers.\par\medskip
\noindent\textit{Editorial note.} rqual sv-mismatch batch fix 2026-10-14 (Chief ruling: re-assemble unpublished only): source\_version aligned verbatim to the source dir license-observation.json sv (package evidence authority); revision bumped for re-assembly. Recipe from lane-s3.\par\medskip
\noindent\textit{Allowed standard prerequisites (editorial).}\begin{itemize}
\item Roelcke colorings, the definable Ramsey property and joint embedding patterns. For a Fraisse structure with Roelcke precompact automorphism group the article studies colorings that are Roelcke, that is, uniformly continuous for the Roelcke uniformity, and among them the colorings attached to a point of the compact space of joint embedding patterns; the age of the structure is the class of its finite substructures, and the definable Ramsey property is the statement that for finite substructures one can find an embedding whose induced coloring is constant on the copies of a smaller substructure, in the sense made precise by the article's definitions of joint embeddings and their patterns. The author's own definitions and the surrounding results are the ones used by the counted statement and proof.. The counted proposition is the bridge the article uses between the dynamical Ramsey property for Roelcke colorings and the combinatorial Ramsey property of the age, and it is the tool with which its main definability theorem is proved.
\end{itemize}\medskip
\ifdefined\hypertarget\hypertarget{sourcepointer-1}{}\fi
\noindent\textit{Source reference.} Original proposition 4.9: sources/ahl-16-fixed-points-2019/body-01.tex, source lines 571--573. Navigation pointer invoked by the retained proof for the extension of Roelcke colorings, which is what makes the definable Ramsey property the finitary form of the dynamical one. That proposition is an earlier result of the article and is not retained in this card, so the reference is declared as a pointer carrying the number the article itself prints for that statement. The article declares its theorem-like environments under a single counter attached to the section, so the printed number is the section and the place of the statement inside it, counted by enumerating every numbered theorem-like environment of the source in order: it is 4.9; the cited statement is an allowed input and none of it is reproduced inside the span.\par\medskip
\ifdefined\hypertarget\hypertarget{sourcepointer-2}{}\fi
\noindent\textit{Source reference.} Original proposition 4.2: sources/ahl-16-fixed-points-2019/body-01.tex, source lines 492--494. Navigation pointer invoked by the retained proof to see that the tree built in the compactness argument is finitely branching, since the automorphism group is Roelcke precompact. That proposition is an earlier result of the article and is not retained in this card, so the reference is declared as a pointer carrying the number the article itself prints for that statement. The article declares its theorem-like environments under a single counter attached to the section, so the printed number is the section and the place of the statement inside it, counted by enumerating every numbered theorem-like environment of the source in order: it is 4.2; the cited statement is an allowed input and none of it is reproduced inside the span.\par\medskip
\begin{prop}\label{prop:DefinableRP}
Let $\mF$ be a Fra\"{\i}ss\'e structure with Roelcke-precompact automorphism group. Then $\mF$ has the Ramsey property for colorings in $\ROB(\Aut(\mF))_{x}$ for every $x\in R(\Aut(\mF))\}$ iff $\Age(\mF)$ has the definable Ramsey property.
\end{prop}

\begin{proof}
From Proposition~\ref{prop:RPExtRoelcke}, it appears that the definable Ramsey property is nothing other than a finitization of the fact that for every $x\in R(\Aut(\mF))$, the Ramsey property holds for colorings in $\ROB(\Aut(\mF))_{x}$. This proves the converse implication.

The direct implication is obtained by a standard compactness argument as follows. Assume that we can find $\mA, \mB, \mZ$ finite substructures of $\mF$ such that for every $\mC\in \Age(\mF)$, there exists a joint embedding $\langle c,z \rangle$ such that no $b\in \binom{\mC}{\mB}$ satisfies that the map $a\mapsto [a,z]$ is constant on $\binom{b(\mB)}{\mA}$. Consider now the sequence $(\mF_{n})_{n\in \N}$ of initial segments of $\mF$ (recall that $\mF$ is based on $\N$ and that $\mF_{n}$ is the substructure of $\mF$ supported by $\{ k : k\leq n\}$). Each comes with some joint embedding pattern $[\phi_{n}, z_{n}]$ witnessing the failure of the definable Ramsey property. Note that we may assume that each $\phi_{n}$ is just the natural inclusion map from $\mF_{n}$ in $\mF$. Closing off this set of joint embedding patterns under initial segments of the first coordinate, we obtain a countable set whose elements are $[\phi_{m}, z_{n}]$, with $m\leq n$. Setting $[\phi_{m}, z_{n}] \leq [\phi_{p}, z_{q}]$ when $m\leq p$ and $[\phi_{m},z_{q}]=[\phi_{m},z_{n}]$, this becomes a countable tree, which is finitely branching since $\Aut(\mF)$ is Roelcke-precompact (Proposition~\ref{prop:Roelcke-precompact}). By K\"onig's lemma, this tree contains an infinite branch, which can be seen as a joint embedding pattern $[\phi, z]$ of $\mF$ and $\mZ$. By construction, there is no $b\in \binom{\phi(\mF)}{\mB}$ such that $a\mapsto [a,z]$ is constant on $\binom{b(\mB)}{\mA}$. Therefore, the Ramsey property for colorings in $\ROB(\Aut(\mF))_{x}$ fails for any $x=[\phi_{1}, \phi_{2}]\in R(\Aut(\mF))$ satisfying $\phi_{1}=\phi$ and $\phi_{2}$ extending $z$ to $\mF$.
\end{proof}

\end{document}
